% TOP_PROBLEM: {"id": 2891, "title": "Kirby Problem 4.15", "category_id": 7, "category": {"id": 7, "name": "topology", "display_name": "Topology", "description": "Properties preserved under continuous deformations.", "slug": "topology", "order_index": 7, "created_at": "2026-07-31T15:26:25.671Z"}, "status": "open"}
% FIRST_POSTED: 2026-10-01
% ATTEMPT_STATUS: unresolved
% !TeX program = lualatex
% Definition and sources only; this file is not an LLM solution attempt.
\documentclass[11pt]{article}
\usepackage[margin=25mm]{geometry}
\usepackage{fontspec}
\setmainfont{Latin Modern Roman}
\usepackage{amsmath,amssymb,mathtools}
\usepackage{unicode-math}
\setmathfont{Latin Modern Math}
\usepackage{xurl,hyperref}
\hypersetup{colorlinks=true,urlcolor=blue}
\setcounter{secnumdepth}{0}
\setlength{\parindent}{0pt}
\setlength{\parskip}{5pt}
\setlength{\emergencystretch}{3em}
\begin{document}
\section*{\texorpdfstring{Kirby Problem 4.15}{Kirby Problem 4.15}}\label{2891}
\subsection{Definitions and mathematical statement}
For a smooth closed oriented four-manifold $M$, the intersection pairing on
$H^2(M,{\mathbb{R}})$ is $Q_M(a,b)=\langle a\smile b,[M]\rangle$. Let $b_2$ be its dimension and $\sigma(M)=b_2^+-b_2^-$ its signature. A spin structure exists precisely when the second Stiefel--Whitney class $w_2(TM)$ vanishes. The conjecture is
\[
 b_2(M)\ge\frac{11}{8}|\sigma(M)|
\]
for every such spin manifold. No simple-connectedness assumption appears in the selected target. Smoothness is essential: the assertion is not a purely algebraic restriction on arbitrary unimodular bilinear forms or an unrestricted claim for topological four-manifolds.
\par\smallskip\textit{Formulation and concept references:} \href{https://web.stanford.edu/~cm5/4D.pdf}{[S1]}.

\subsection{Short English statement}
Must the second Betti number of every smooth closed spin four-manifold be at least eleven-eighths the absolute value of its signature?
\subsection{Research attempt: closure operations and the intersection-form deficit}
\textbf{Status.} We prove preservation results and compute the exact numerical deficit on a standard family of forms. These do not prove the conjecture for every smooth spin four-manifold. The smoothness hypothesis is retained throughout. The primary survey [S1] supplies the external facts about the K3 surface and Furuta's theorem used below.

\subsubsection{1. Connected sums preserve the desired inequality}
For closed oriented connected four-manifolds, Mayer--Vietoris for removal and gluing of balls identifies the real intersection pairing of $M\#N$ with the orthogonal sum of those of $M$ and $N$. Therefore
\[ b_2(M\#N)=b_2(M)+b_2(N),\qquad \sigma(M\#N)=\sigma(M)+\sigma(N).\]
Spin structures on the punctured manifolds extend across the neck: the sphere $S^3$ has a unique spin structure, so the restrictions can be identified. If both summands obey the target inequality, the triangle inequality gives
\[b_2(M\#N)\ge\tfrac{11}{8}(|\sigma(M)|+|\sigma(N)|)\ge\tfrac{11}{8}|\sigma(M\#N)|.\]
Orientation reversal changes the sign of signature and leaves $b_2$ unchanged, so it preserves the assertion as well. This argument also shows why cancellation of signatures makes the desired bound easier, not harder.

\subsubsection{2. Concrete spin families}
Every closed oriented surface $\Sigma_g$ is spin because $w_2(T\Sigma_g)$ is its Euler class modulo two and $2-2g$ is even. The product is spin by the Whitney formula. K\"unneth gives
\[b_2(\Sigma_g\times\Sigma_h)=2+4gh.\]
An orientation-reversing diffeomorphism of the first surface, times the identity of the second, reverses the product orientation. Naturality of the intersection pairing then gives $\sigma=-\sigma$, hence signature zero. Thus every such product satisfies the inequality without requiring a simply connected hypothesis.

As an external, standard geometric example, a K3 surface is spin with $b_2=22$ and $\sigma=-16$; it realizes equality. Its orientation reverse does too. Connected sums of $r$ copies with the same orientation satisfy $b_2=22r$ and $|\sigma|=16r$, again equality. Adding $S^2\times S^2$ contributes two to $b_2$ and zero to signature. No statement that every spin four-manifold decomposes into these examples is used.

\subsubsection{3. The numerical gap in even indefinite forms}
Consider a form consisting of $2a$ copies of a definite $E_8$ form, all of the same sign, and $b$ hyperbolic planes $H$, where $a,b\ge0$. Its rank and absolute signature are
\[r=16a+2b,\qquad |s|=16a.\]
The target inequality is exactly $b\ge3a$. On a smooth spin manifold with indefinite form, the external $10/8$ theorem quoted in [S1] instead gives $r\ge(10/8)|s|+2$, hence $b\ge2a+1$. For $a=1$ these coincide; for $a\ge2$ the interval $2a+1\le b<3a$ is not eliminated by that estimate alone. This is an exact accounting of the missing strength, rather than an interpolation between the two coefficients.

The algebraic form $2E_8$ has rank 16 and signature 16 and violates $11/8$. It is even and unimodular. This is only an algebraic warning: we have not produced a smooth spin manifold realizing it. In particular, a construction of an intersection matrix without a smooth realization cannot refute the question.

\subsubsection{4. Adversarial audit and remaining gap}
The connected-sum argument uses absolute values in the correct direction; a lower bound on a sum of absolute signatures would not follow from the absolute signature of the sum. The use of the $10/8$ theorem above is restricted to the indefinite case in the cited formulation. The product calculation and closure results establish genuine families, but no surgery, decomposition, or gauge-theoretic argument here upgrades $2a+1$ to $3a$ for general smooth spin forms. That upgrade, together with the general geometric cases not encoded by these examples, remains unresolved.

\subsection{Sources}
\begin{itemize}
\item[S1]Ciprian Manolescu, Four-dimensional topology, Section 4.4 (2024).
\url{https://web.stanford.edu/~cm5/4D.pdf}.
\textit{Formulation source.}
\end{itemize}
\end{document}
